% Chapter 1: Introduction
\section{Introdução}

A introdução serve para apresentar o tema do trabalho, contextualizando o leitor sobre a área de estudo. Aqui, descrevemos a importância do ensino de química, especificamente da Tabela Periódica, e introduzimos o conceito de jogos sérios como uma abordagem pedagógica inovadora.

Este trabalho propõe o desenvolvimento de "Alchemon", um jogo sério projetado para facilitar a aprendizagem dos elementos químicos. A motivação para este projeto vem da observação de que métodos de ensino tradicionais podem não ser suficientemente engajadores para a geração atual de estudantes, que estão imersos em mídias digitais. Como aponta Gee em seu trabalho sobre video games e aprendizagem, ambientes de jogo bem projetados podem ser máquinas de aprendizado poderosas \cite{journal_example}.

O objetivo é criar uma ferramenta que seja ao mesmo tempo educativa e divertida, transformando o estudo da Tabela Periódica em uma experiência interativa e memorável. Acreditamos que a gamificação pode aumentar o interesse dos alunos e melhorar a retenção do conteúdo, conforme discutido em diversas publicações da área.

\begin{figure}[h]
    \centering
    % \includegraphics[width=0.8\columnwidth]{placeholder} % Descomente e adicione sua imagem
    \caption{Exemplo de um diagrama de fluxo do jogo. Uma imagem pode ser inserida aqui para ilustrar um conceito. A pasta padrão para imagens é 'figuras/'.}
    \label{fig:fluxo_jogo}
\end{figure}
